% !TEX root = ../../main.tex
%% Corrigé du défi du chapitre 8 — vérifié par verifications/defis_06_10.py.

\begin{blocLivre}{cor-8-25a}
\begin{corrige}{25}{Le groupe du carré et le comptage des coloriages}
\demarche{Tout calcul dans $D_4$ se ramène aux deux règles $r^{4}=s^{2}=\mathrm{id}$
et $sr=r^{-1}s$ ; un coloriage invariant est constant sur chaque cycle.}
\partieExo{A}{Les huit isométries}
\begin{questions}
  \item L'image de $A_1$ est l'un des $4$ sommets ; celle de $A_2$ est un
    sommet à distance $A_1A_2$ de celle de $A_1$, donc l'un de ses $2$
    voisins : au plus $4\times2=\res{8}$ éléments. En permutations :
    $r=(1\ 2\ 3\ 4)$ et $s=(2\ 4)$.
  \item $srs(1)=s(r(1))=s(2)=4$, $srs(2)=s(r(4))=s(1)=1$,
    $srs(3)=s(r(3))=s(4)=2$, $srs(4)=s(r(2))=s(3)=3$ : ainsi
    $srs=(1\ 4\ 3\ 2)=r^{3}=\res{r^{-1}}$. Les $r^{k}$ ($0\leqslant k\leqslant3$)
    sont distincts et préservent l'orientation, les $r^{k}s$ la renversent
    et sont distincts entre eux : on a $8$ éléments distincts, donc tout
    $D_4$. Il n'est \res{\text{pas commutatif}} : $sr=r^{-1}s=r^{3}s\neq rs$.
  \item $rs=(1\ 2)(3\ 4)$ : réflexion d'axe la médiatrice de
    $[A_1A_2]$ ; $r^{2}s=(1\ 3)$ : axe $(A_2A_4)$ ;
    $r^{3}s=(1\ 4)(2\ 3)$ : médiatrice de $[A_1A_4]$. Ordres :
    $\mathrm{id}$ : $1$ ; $r$, $r^{3}$ : $\res{4}$ ; $r^{2}$ et les
    quatre réflexions : $\res{2}$.
\end{questions}
\partieExo{B}{Structure du groupe}
\begin{questions}[resume]
  \item $\langle r\rangle$ est cyclique d'ordre $4$. Pour
    $V_1=\ens{\mathrm{id},r^{2},s,r^{2}s}$ : $sr^{2}s=(srs)^{2}=r^{-2}=r^{2}$,
    d'où $s\cdot r^{2}s=r^{2}$, $r^{2}\cdot s=r^{2}s$,
    $(r^{2}s)^{2}=r^{2}(sr^{2}s)=r^{4}=\mathrm{id}$ : stable, et chaque
    élément est son propre inverse. Même calcul pour
    $V_2=\ens{\mathrm{id},r^{2},rs,r^{3}s}$. Dans $V_1$ et $V_2$, tout
    élément non trivial est d'ordre $2$ : ils \res{\text{ne sont pas
    cycliques}} (groupes de Klein), seul $\langle r\rangle$ l'est.
  \item $r^{k}$ commute avec $s$ si et seulement si
    $r^{k}s=sr^{k}=r^{-k}s$, i.e. $r^{2k}=\mathrm{id}$, i.e. $k\in\ens{0,2}$.
    Une réflexion $r^{k}s$ ne commute pas avec $r$ :
    $(r^{k}s)r=r^{k-1}s\neq r^{k+1}s=r(r^{k}s)$. Enfin $r^{2}$ commute
    avec $r$ et avec $s$, donc avec tout : le centre est
    $\res{\ens{\mathrm{id},r^{2}}}$.
  \item Un isomorphisme transporte la commutativité, et $D_4$ n'est pas
    commutatif. Le groupe commutatif $\Z/4\Z\times\Z/2\Z$ est d'ordre $8$
    et $(\overline1,\overline0)$ y est d'ordre $4$ : avoir les mêmes
    ordres d'éléments ne suffit pas à être isomorphes.
\end{questions}
\end{corrige}
\end{blocLivre}

\begin{blocLivre}{cor-8-25b}
\begin{corrigeSuite}
\partieExo{C}{Compter les coloriages à symétrie près}
\begin{questions}[start=7]
  \item Un coloriage $c$ est invariant par $g$ si et seulement si
    $c(g(i))=c(i)$ pour tout $i$, c'est-à-dire si $c$ est constant sur
    chaque cycle de $g$ : on choisit librement une couleur par cycle,
    d'où $\mathrm{Fix}(g)=2^{c(g)}$.
    \begin{center}
    \begin{tabular}{@{}lcccccccc@{}}
      $g$ & $\mathrm{id}$ & $r$ & $r^{2}$ & $r^{3}$ & $s$ & $rs$ & $r^{2}s$ & $r^{3}s$\\\hline
      $c(g)$ & $4$ & $1$ & $2$ & $1$ & $3$ & $2$ & $3$ & $2$\\
      $\mathrm{Fix}(g)$ & $16$ & $2$ & $4$ & $2$ & $8$ & $4$ & $8$ & $4$
    \end{tabular}
    \end{center}
  \item Somme : $48$, d'où $\dfrac{48}{8}=\res{6}$ classes. Directement,
    une classe est déterminée par le nombre de sommets noirs ($0$ à $4$),
    sauf pour deux sommets noirs, qui peuvent être voisins ou opposés :
    $5+1=6$ classes (figure).
  \item Avec trois couleurs, $\mathrm{Fix}(g)=3^{c(g)}$ :
    $\dfrac{81+3+9+3+27+9+27+9}{8}=\dfrac{168}{8}=\res{21}$ colliers.
\end{questions}
\begin{center}
\begin{tikzpicture}[line cap=round]
  \foreach \cfg [count=\i] in {{0,0,0,0},{1,0,0,0},{1,1,0,0},{1,0,1,0},{1,1,1,0},{1,1,1,1}}{%
    \begin{scope}[xshift={(\i-1)*2.2cm}]
      \draw[otNavy,line width=0.7pt] (0.6,0) -- (0,0.6) -- (-0.6,0) -- (0,-0.6) -- cycle;
      \foreach \b [count=\k] in \cfg {%
        \pgfmathsetmacro{\ang}{(\k-1)*90}
        \ifnum\b=1
          \fill[otInk] (\ang:0.6) circle (2.6pt);
        \else
          \filldraw[fill=white,draw=otInk,line width=0.6pt] (\ang:0.6) circle (2.6pt);
        \fi}
    \end{scope}}
\end{tikzpicture}\\[0.5mm]
\legende{Les six coloriages à symétrie près : $0$, $1$, $2$ voisins,
$2$ opposés, $3$ et $4$ sommets noirs.}
\end{center}
\end{corrigeSuite}
\end{blocLivre}
